%=======================================================
\chapter{Tabelas}

\textsc{Quais os pacotes?} \textbf{tabularx} (Acesso: \url{https://www.ctan.org/pkg/tabularx}) e \textbf{booktabs} (Acesso: \url{https://www.ctan.org/pkg/booktabs/}).

\textsc{O que fazem?}: Ambos permitem maior controle de personalização nas tabelas. 

Relevante observar que tabelas, \textit{stricto sensu}, são elementos destinados a apresentar dados \textbf{numéricos} e não apresentam linhas laterais. Temos o costume, aqui no Brasil, de incluir tais linhas, inserir outros dados que não sejam numéricos e identificar o elemento como tabela. Esse é um procedimento incorreto. 

Caso se apresentem dados não numéricos e linhas laterais, deve-se identificar esse elementos como \textsc{quadro}.

Veja a seguir um exemplo de como normalmente se faz no Brasil\footnote{Obviamente que fica a critério da pessoa que elabora o trabalho optar pela forma abaixo ou não.}.
\ \\

\begin{codex}{Código para tabela com legenda e referência cruzada}
\begin{table}[h!]
	\begin{center}
	\begin{tabularx}{0.8\textwidth} { 
	| >{\centering\arraybackslash}X 
	| >{\centering\arraybackslash}X 
	| >{\centering\arraybackslash}X | }
	\hline
	item 1 & item 2 & item 3 \\
	\hline
	item 4  & item 5  & item 6  \\
	\hline
	\end{tabularx}
	\caption{Tabela de teste com linhas laterais.}
	\label{table:1}
	\end{center}
	\end{table}
\end{codex}
\ \\

\textsc{Resultado:}

\begin{table}[h!]
	\begin{center}
		\begin{tabularx}{0.8\textwidth} { 
				| >{\centering\arraybackslash}X 
				| >{\centering\arraybackslash}X 
				| >{\centering\arraybackslash}X | }
			\hline
			item 1 & item 2 & item 3 \\
			\hline
			item 4  & item 5  & item 6  \\
			\hline
		\end{tabularx}
		\caption{Tabela de teste com linhas laterais.}
		\label{table:1}
	\end{center}
\end{table}

E para referenciar a tabela/quadro, usamos o comando \verb|\ref{table:1}|, como abaixo.

Ex.: \textit{Como podemos observar na tabela \ref{table:1}, blá blá...}

No entanto, a correta aplicação e formatação de uma tabela deve ser conforme o código a seguir.
\ \\

\begin{codex}{Código de exemplo de tabela}
\begin{table}[h!]
	\begin{center}
	\begin{tabular}{c|c|c|c}
	\hline
	Percentual & 15\% & 37\% & 50\% \\
	\hline
	Número & 23 & 44 & 71 \\
	\hline
	Variação & 2.78 & 4.59 & 10.96 \\
	\hline
	Total & > 3 & < 8 & >15  \\
	\hline
	\end{tabular}
	\caption{Tabela de teste sem linhas laterais. Abaixo, a ref. cruzada.}
	\label{table:2}
	\end{center}
\end{table}
\end{codex}
\ \\

\textsc{Resultado:}

\begin{table}[h!]
	\begin{center}
		\begin{tabular}{c|c|c|c}
			\hline
			Percentual & 15\% & 37\% & 50\% \\
			\hline
			Número & 23 & 44 & 71 \\
			\hline
			Variação & 2.78 & 4.59 & 10.96 \\
			\hline
			Total & > 3 & < 8 & >15  \\
			\hline
		\end{tabular}
		\caption{Tabela de teste sem linhas laterais. Abaixo, a ref. cruzada.}
		\label{table:2}
	\end{center}
\end{table}

E, como acima, referenciamos a tabela (ref. cruzada) com o comando \verb|\ref{table:2}|.

Ex.: \textit{Como podemos observar na tabela \ref{table:2}, blá blá...}

A lista de tabela, quando identificadas pelos dois comandos \verb|\caption{...}| e \verb|\label{...}| já são inseridas automaticamente na seção \textbf{Lista de tabelas}.

Uma observação importante é que a numeração da tabela é feita conforme a seção ou subseção, inicialmente, seguida do número (se primeira tabela, se segunda, etc) daquela seção. No exemplo acima temos tabela \ref{table:2} porque ela é a segunda tabela da seção 9.

Tabelas em \LaTeX\ são incrivelmente \qq{complicadas} de fazer. Sugiro abaixo uma breve lista com geradores de códigos de tabela em \LaTeX\ \textit{on-line}.

\begin{enumerate}
	\item \url{https://www.tablesgenerator.com/}
	\item \url{https://www.latex-tables.com/}
	\item \url{https://clevert.com.br/latex/latextable.php}
\end{enumerate}

Por fim, ao final do arquivo \verb|pretextuais.tex|, há os seguintes códigos:
\ \\
\begin{codex}{Códigos para tabelas e figuras}
	\newpage\tableofcontents*
	\thispagestyle{empty}
	\newpage\listoffigures*
	\thispagestyle{empty}
	\newpage\listoftables*
	\thispagestyle{empty}
\end{codex}
\ \\
Caso não existam nem figuras nem tabelas no trabalho, basta excluir as linhas: 

\begin{verbatim} 
	\newpage\listoffigures*
	\thispagestyle{empty}
	\newpage\listoftables*
	\thispagestyle{empty} 
\end{verbatim}